\cotefile{149}

\scanpage{79}
\mspage{39}
l'hom ``local''
\diagram{(12)\quad $\unsure{\kappa|k}_{i_0}:$\quad
\begin{tikzcd}[ampersand replacement=\&,row sep=tiny]
T \arrow[r] \arrow[d,equal] \& \pi_1(U_{\unsure{\eta|k}}) \\
T_\infty(k) \&
\end{tikzcd}}
associé à l'indice $i_0$. Je vais admettre
ce point, qui ne peut guère être difficile.

\margin{\deleted{* En fait, on n'a pas $=$ besoin
que le centre de $\pi$ soit trivial
pour cette construction --- c'est un
fait général pour les situations fibrées,
on trouve une ext. du $\pi_1$ de la
base par le $\pi_1$ de la fibre, \unsure{comme}}}

Pour terminer ce n°, je veux encore
étudier, dans la situation d'une
$U$ \deleted{élémentaire}\inserted{courbe} relative sur un $S$ avec
\deleted{\unsure{du}} $U=X\setminus T$, \inserted{$X$ lisse et propre sur $S$,} $T$ fini étale, une section
$g_i$ donnée de $T$ sur $S$, les ``sections de
2ième espèce'' de l'extension
\[
1-\pi-E_U\longrightarrow E_S-1
\]
associées à $i=i_0$ --- qui définit une \inserted{classe de $\pi$-conj. de} \unsure{ss-groupe|ss-)groupe}
\unsure{fermés|libre} $L_{i_0}$ dans $\pi$. (On suppose qu'on
a bien une telle suite exacte, i.e. que
$\pi_2(S)\longrightarrow\pi_1(\text{fibre})$ est nul, ce qui sera
le cas si $\pi_2(S)=0$, p.ex. (je présume)

\margin{en tous cas,
$=$ sans $\pi_1(S)\longrightarrow\pi$ trivial,
$E_S=\pi_1(S)$ opère extérieurement
sur $\pi=\pi_1(\text{fibre géom})$, d'où
une extension puisque $\pi$
est à centre trivial... Cette extension
admet alors * $\pi_1(U)$ comme
quotient...}
\margin{il faudrait
quand même
\unsure{les|le} vérifier !}

Si on est dans le cas d'un modèle élémen-
taire au dessus d'un corps de car.
zéro (la restriction de car. nulle
étant sans doute inutile, si on se borne
\margin{78}

\scanpage{80}
aux groupes fondamentaux premiers aux car résiduelles\dots) Par un argument général, le normalisateur de $L_i$ dans $E_U$ s'envoie \uline{sur} $E_S$, on trouve donc des scindages \unsure{pour|par} cette extension, qu'on peut regarder comme une extension
\[
(13)\quad 1\unsure{\longrightarrow|-}T\longrightarrow N(L_i)\longrightarrow E_S\longrightarrow1
\]
La classe d'isomorphie est un élément \deleted{\unsure{de}}
\[
(14)\quad c_i\in H^2(E_S,\underset{\substack{\Vert\\\varprojlim\mu_n}}{T}),
\]
que je me propose d'étudier. On a, si $S$ est un $K(\pi,1)$
\[
(15)\quad H^2(E_S,T)\simeq H^2(S,T)\simeq\varprojlim_n H^2(S,\mu_n)
\]
d'ailleurs on a une suite exacte de Kummer (\unsure{où} $\operatorname{Pic}(S)=H^1(S,G_m)$)
\[
(16)\quad 0\longrightarrow\operatorname{Pic}(S)_n\unsure{\longrightarrow|-} H^2(S,\mu_n)\longrightarrow {}_nH^2(S,G_m)\longrightarrow0
\]
d'où par passage à la limite
\[
(17)\quad 0\unsure{-|\longrightarrow}
\underset{\substack{\Vert\,\text{déf}\!^{\mathrm n}\\\varprojlim(\operatorname{Pic}_n=\operatorname{Pic}/n\operatorname{Pic})}}{\operatorname{Pic}(S)^\wedge}
\longrightarrow H^2(S,T)\longrightarrow\varprojlim {}_nH^2(S,G_m)
\]
Dans le cas où $S$ est un schéma. élémentaire sur un corps de \deleted{\unsure{car}}. \unsure{abst.|de t.}
\margin{79}

\scanpage{81}
\mspage{40}
de type fini sur $Q$, $\operatorname{Pic}(S)$ est un $\mathbb Z$-\deleted{\unsure{groupe}}\inserted{module} de type fini \inserted{(par Mordell-Weil-Néron)}, donc l'hom
\[
(18)\quad \operatorname{Pic}(S)\longrightarrow\operatorname{Pic}(S)^\wedge\hookrightarrow H^2(S,T)
\]
est \uline{injectif}.

Sans nous soucier de cette condition, \deleted{\unsure{pour}} considérons le cas général --- je dis que la classe \deleted{\unsure{de}} $c$ \inserted{(14)} est dans l'image de (18), de façon précise que c'est l'image de l'élément
\[
(19)\quad \gamma_i\in\operatorname{Pic}(S)
\]
classe du faisceau conormal (ou normal, le diable sait) de $X$ le long de $g_i$. Principe d'une vérification : on prend le complété formel de $X$ le long de $g_i(S)$ et \deleted{\unsure{le}} \unsure{qui ouvre|qu'on ouvre} --- on enlève $g_i(S)$, et on interprète la suite exacte (13) comme la suite exacte d'homotopie de ce \deleted{\unsure{schéma}} topos \unsure{monté|inventé}, au dessus de \deleted{\unsure{Spec}} $S$. On a donc à prouver une histoire d'ombres\dots
\margin{80}

\scanpage{82}
Dans le cas ``arithmétique'', on voit donc que l'ext. (13) est scindée sss $\unsure{\gamma_i|\xi_i}=0$ i.e. \unsure{si l'on|ssi on} peut trouver, globalement sur $S$, une uniformisante sur $X$ le long de $g(S)$; et cette condition (tjrs satisfaite loc.\inserted{(Zar)} sur $S$) est en tous cas tjrs suffisante (cas arith. ou pas !)

Quand $\gamma_i=0$, parmi les scindages, il y a ceux provenant du choix d'une \deleted{section}\inserted{base} de $J_i/J_i^2$ qui sont privilégiés L'indétermination du choix d'une telle base est dans $G_m(S)=H^0(S,G_m)$ --- celle du choix d'une section de (13) est dans
\[
(20)\quad H^1(S,T)\simeq\varprojlim H^1(S,\mu_n)
\]
On a ici des suites exactes de Kummer
\[
\deleted{\unsure{(21)}}\quad 0\longrightarrow\deleted{\unsure{H^0(S)_n}}G_m(S)_n\longrightarrow H^1(S,\mu_n)\longrightarrow{}_n\operatorname{Pic}(S)\longrightarrow1
\]
d'où par passage à la limite
\[
(21)\quad 0-G_m(S)^\wedge\longrightarrow H^1(S,T)\longrightarrow\varprojlim_n{}_n\operatorname{Pic}S\unsure{?}
\]
Dans le ``cas arithmétique'' le troisième terme est nul, on trouve donc
\[
\deleted{\unsure{G_m}}\quad H^1(S,T)\simeq G_m(S)^\wedge
\]
\margin{81}

\scanpage{83}
\mspage{41}
Si le genre est zéro, prenant une de ces sections de $T$ sur $S$ comme section à l'infini, OPS (quitte à se localiser) $U_S=\mathbb E'_S\smallsetminus T'$, donc $f$ s'identifie à une section de $\mathbb E^1_S$ sur $S'$, i.e. de $\underline O_S$ sur $S'$, donc (comme $\operatorname{codim}(Z,S)\geq2\deleted{\unsure{1}}$, $S'$ normal) elle se prolonge en une section de $\mathbb E'_S$. Et on achève comme précédemment, OK

[Considérons \inserted{donc} les diviseurs irréductibles $D_i$ sur $S$, ou ce qui revient au $=$, les pts $x_i$ de $X$ de codim $1$, i.e. tels que $\underline O_{x_i}$ soit un anneau-jauge (anneau de valuation discrète). Considérons son normalisé $\overline{\underline O}_{x_i}$ dans $\overline K$, choisissons un idéal maximal $\widetilde x_i$ de $\overline{\underline O}_x$ \deleted{\unsure{ce|et} qui \unsure{donne|devra|leur}}. (ces idéaux correspondent aux pts du normalisé $\widetilde S$ de $S$ dans $\overline K$ \inserted{au dessus de $x$}) et considérons son stabilisateur $N_x$ dans \unsure{le|$E$} $E_K$, qui opère \inserted{donc} dans $k(\widetilde x)=\overline{k(x)}$, et s'envoie en fait, on le sait, \uline{sur} $\operatorname{Gal}(\overline{k(x)}/k(x))$
\margin{82}

\scanpage{84}
Soit $I_{\widetilde x}$ le noyau de l'hom. obtenu
(ou stabilisateur ``géométrique'' de $\widetilde x$), donc
on a une suite exacte
\[
(7)\quad 1-I_{\widetilde x}\longrightarrow N_{\widetilde x}\longrightarrow\operatorname{Gal}(\overline{k(\unsure{x|\dot{x}})}/k(x))\longrightarrow1
\]
et par Kummer un iso. can.
\[
(8)\quad I_{\widetilde x}\simeq T_\infty(\overline{k(x)}^*)\simeq T_\infty(\overline K^*)
\]
\margin{à cause de
l'hyp de
car. 0 !}
On notera que si $x$ est le pt gén. du
diviseur $D$, alors $k(x)$ est le corps des
fonctions de $D$. C'est un corps de
type fini sur $\mathbb Q$.

Il est immédiat (sans supposer que le
corps de base pour $S$ soit $\mathbb Q$) que
le noyau de $\deleted{\unsure{\pi|\Gamma}}\inserted{E}_K\longrightarrow E_S$ est le sous-groupe
invariant engendré par les $I_{\widetilde x}$. Donc
l'hypothèse que $E_K\longrightarrow E_{U_S}$ s'annule
sur le dit noyau, signifie aussi qu'il
s'annule sur chacun des $I_{\widetilde x}$.

Soit \unsure{donc|sur} $U_{\underline O_x}$ induit par $U$ sur
$\operatorname{Spec}O_x$, on a donc des factorisations
\margin{83}

\scanpage{85}
\mspage{42}
d'ailleurs $G_m(S)$ n'a pas d'él. infiniment divisible\deleted{\unsure{s}}, donc
\[
\unsure{(22)|(21)}\quad G_m(S)\longrightarrow G_m(S)^\wedge\simeq H^1(S,T)
\]
\margin{(cas ``arithm'')}
est injectif. Ainsi, quand $\gamma_i=0$ i.e. quand (13) admet des scindages ``\unsure{géométriques|géométrique}' \inserted{(\unsure{car} il suffit pour cela qu'il existe une section tt court)} ceux-ci forment un torseur sous $\deleted{\unsure{H^1(S,T)}\unsure{=|\simeq}}G_m(S)^{\unsure{\wedge}}$, qui s'identifie à un sous-torseur du $(H^1(S,T)=G_m(S)^\wedge)$-torseur de tous les scindages de (13). Pour que la ``description profinie de la géom. alg. absolue sur $\mathbb Q$'' soit complete, il y faudrait également caractériser (en termes de \inserted{cette} description profinie) ce sous-ensemble remarquable.

Je voudrais enfin comprendre encore comment une section d'\unsure{extension|extensions} du type (1) peut se ``spécialiser'' en une section du type (2), dans le cas de \unsure{corps de} valuation. Pour ceci, je reprends
\margin{84}

\scanpage{86}
la sempiternelle situation de $U\unsure{=|:}X\setminus T$ sur $S$, avec une $g_i=g_{i_0}$ section de $T$ sur $I$, et je suppose \unsure{donné|donnée} une section rationnelle $f$ de $U$ sur $S$, qui en un point particulier $s$ de \deleted{\unsure{car.}} codim $1$ va coïncider avec $g_{i_0}$ en tt que section de $X$
\[
g_i(s)=f(s)
\]
de sorte que l'hom induit sur le groupe d'inertie $I_{s'}\subset E_{\deleted{\unsure{S|s}}\inserted{K}}$ par le composé $E_K\longrightarrow E_{U_K}-E_U$ (qui envoie en fait $I_{s'}$ dans $\pi$, et $=$ dans les $L_i$, on l'a vu) soit non \deleted{nulle} triviale. Considérons alors le groupe de décomposition $N_{s'}$ (cf suite exacte (7), dont on a changé les notations), qui s'identifie comme le $\pi_1$ de la ``rondelle trouée'' correspondant au hensélisé (pas strict) de la jauge $\underline O_{S,s}$ --- ou mieux, le $\pi_1$ du fibré en
\margin{85}

\scanpage{87}
\mspage{43}
disques troués autour du germe de
variété $D=(\overline s)_{\mathrm{red}}$ en $s\dots$, et l'extension
de $N_{s'}$ par $\pi$ contenue dans l'ext. $E_{U_K}$ de
$E_K$ par $\pi_1$.
\margin{\deleted{\unsure{c'est l'extension} $E_{U_{\unsure{\widetilde K|K}}}$ \unsure{correspondante|correspondant}} \unsure{\deleted{\unsure{à le $s$|à la $s$}}}}
La section donnée (plutôt,
\unsure{construit|construite} via $f_K$) de l'ext. $E_{U_K}$, en
définit une de la sous extension. \deleted{\unsure{Si on|Comme}}

L'interprétation géométrique de ces constructions
\deleted{\unsure{considérer|considère} d'autre part les scindages \unsure{de|des}}
profinies est \unsure{la|le} \unsure{suivante|suivant} : Soit
$\deleted{\unsure{(\widetilde S)_w|(\widetilde S)_s|\widetilde Z}}\inserted{\unsure{\widetilde{\underline O}_s|\underline O_s}}$ le hensélisé (pas strict) de $\underline O_{S,s}=\underline O_s$,
et $\widetilde S=\operatorname{Spec}(\widetilde{\underline O}_s)$ \inserted{(trait hensélien)}, $\widetilde K$ son corps des
fonctions, on a $U_{\widetilde S}$ sur $\widetilde S$, $U_{\widetilde K}$ sur $\widetilde K$,
et \deleted{l'extension}
\[
(23)\quad N_{s'}\quad
\begin{gathered}[t]
\text{(groupe de}\\
\text{décomposition}\\
\text{de }E_{K,\overline K}\text{ en }s')
\end{gathered}
\simeq E_{\widetilde K}
\]
et \unsure{la|le} \deleted{\unsure{suite|seule|sous}}\inserted{$E$} extension de ce groupe par $\pi$
n'est autre que
\diagram{\begin{tikzcd}[ampersand replacement=\&,column sep=large,row sep=1.2em]
(24) \& 1 \arrow[r,dash] \& \pi \arrow[r] \& E_{U_{\widetilde K}} \arrow[r] \arrow[ddd] \& E_{\widetilde K} \arrow[r] \arrow[ddd,bend right=18] \& 1 \\
\& \& \begin{gathered}\Vert\\\pi_1(U_{\overline K})\end{gathered} \& \& \begin{gathered}\Vert\\\operatorname{Gal}(\overline K/\widetilde K)\end{gathered} \& \\
\text{s'envoyant dans} \& \& \& \& \& \\
(25) \& 1 \arrow[r,dash] \& \pi \arrow[r] \& E_U \arrow[r] \& E_{\widetilde O} \arrow[r] \& 1 \\
\& \& \& \& \begin{gathered}\text{\tikz[baseline=(rel.base)]{\node[rotate=90,inner sep=0pt] (rel) {$\simeq$};}}\\\operatorname{Gal}(\overline k/k)\end{gathered} \&
\end{tikzcd}}
($k=k(s)$ est le corps résiduel).
\margin{86}

\scanpage{88}
On trouve donc un diagramme de suites exactes
\diagram{\begin{tikzcd}[ampersand replacement=\&,column sep=large,row sep=large]
 \& \& \& 1 \arrow[d] \& 1 \arrow[d] \& \\
 \& \& \& \deleted{I_{\unsure{s'}}}\inserted{J} \arrow[r,phantom,"\simeq"] \arrow[d] \& I_{s'}\simeq T \arrow[d] \& \\
 \& 1 \arrow[r] \& \pi \arrow[r] \arrow[d,equal] \& E_{U_{\widetilde K}} \arrow[r,dash] \arrow[d] \& \underset{\substack{\tikz[baseline=-0.5ex]{\node[inner sep=0pt,rotate=120] {$\scriptstyle\simeq$};}\\\operatorname{Gal}(\overline K/\widetilde K)}}{E_{\widetilde K}} \arrow[r] \arrow[d] \& 1 \\
(26) \& 1 \arrow[r,dash] \& \pi \arrow[r,dash] \& E_{U_{\widetilde O}} \arrow[r,dash] \arrow[d] \& \underset{\substack{\tikz[baseline=-0.5ex]{\node[inner sep=0pt,rotate=-60] {$\scriptstyle\simeq$};}\\\operatorname{Gal}(\overline k/k)}}{E_{\widetilde O}} \arrow[r] \arrow[d] \& 1 \\
 \& \& \& 1 \& 1 \&
\end{tikzcd}}
le noyau \inserted{$J$} de $E_{U_{\widetilde K}}-E_{U_{\widetilde O}}$ s'envoyant isomorphiquement par $E_{U_{\widetilde K}}\longrightarrow E_{\widetilde K}$ sur le sous-groupe d'inertie $I_{s'}$, lui-$=$ can. isom. à $T$.

La section de \inserted{$E_{U_{\widetilde K}}$ sur} $E_{\widetilde K}\simeq N_{s'}$, induite par la section de départ de $E_{U_K}$ sur $E_K$ (elle $=$ \unsure{définie|definie} par \inserted{le pt $K$-rat.} $f_K$ de $U_K$ sur $K$) n'est autre bien sûr que celle définie par le pt $\widetilde K$-rationnel correspondant $f_{\widetilde K}$ de $U_{\widetilde K}$. Son composé avec $E_{U_{\widetilde K}}-E_{U_{\widetilde O}}$ envoie $I\simeq T$ dans $\pi$, on a vu comment : essentiellement par un $\unsure{\kappa_i|k_i}\circ(\mu.\operatorname{id}_T)$, où $\unsure{\kappa_i|k_i}:T-\pi$ est un hom.-\unsure{lacets|lacet}, relatif à l'indice $i$.

Ainsi \inserted{le groupe de décomposition} $E_{\widetilde K}\deleted{\unsure{(\widetilde f_K)|(\widetilde j_K)|(s')}}$ s'envoie dans le normalisateur de $L_i=\unsure{\kappa_i|k_i}(T)$, dans $E_{U_{\widetilde O}}$. D'ailleurs
\margin{87}

\scanpage{89}
\mspage{44}
la dernière ligne dans (26) s'identifie aussi à
\[
(27)\qquad 1\unsure{-|\longrightarrow}\pi\unsure{-|\longrightarrow}E_{U_s}\longrightarrow E_{k(s)}\unsure{-|\longrightarrow}1
\]
associé à $U_s$, fibre de $U/S$ en $s$, sur le corps résiduel $k(s)$.

\uline{Résumons} La situation de la courbe relative $U/S$, et d'un $s\in S$ de codim $1$, \unsure{définit|definit} par hensélisation (\inserted{\uline{pas}} stricte) en $s$ le diagramme (2\deleted{\unsure{6|5}}\inserted{\unsure{5|6}}), dont la dernière ligne est can. isom. à (27) (par \deleted{\unsure{l'inclusion}}\inserted{l'hom} \inserted{sur \unsure{les|le} $\pi_1$} correspondant :
\diagram{\begin{tikzcd}[ampersand replacement=\&]
U_s \arrow[r,no head] \& U_{\widetilde O} \arrow[dd] \\
\unsure{\vert|\downarrow} \& \\
\unsure{s|k} \arrow[r] \& \operatorname{Spec}\widetilde O
\end{tikzcd}}
)

Une section rationnelle \inserted{$f_{\deleted{\unsure{K}}}$} de $U$ sur $S$, plus gén. \inserted{$\widetilde f$} de $U_{\widetilde O}$ sur $\widetilde S$, définit une section \inserted{$E_{\widetilde f}$} de la ligne médiane de (26), ou ce qui revient au $=$ un relèvement
\diagram{\begin{tikzcd}[ampersand replacement=\&,column sep=large]
(28) \& E_{\widetilde K} \arrow[r,"\lambda"] \& E_{U_{\widetilde O}} \\
 \& \& E_{U_s} \arrow[u,"\sim"']
\end{tikzcd}}
de $E_{\widetilde K}\longrightarrow E_{\widetilde O}$ ($\xleftarrow{\sim}E_{k(s)}$) --- où $E_{\widetilde K}$ \unsure{peut|par} d'ailleurs \unsure{s'identifier|s'identifie} au groupe de décomposition de $\unsure{\underline O_{S,s}|\underline U_{S,s}}$ dans $\overline K$. Les conditions suivantes sont
\margin{88}

\scanpage{90}
équivalentes

a) Cet hom est injectif

b) La restriction de cet hom : $I_{s'}\simeq T$ est
de la forme $\unsure{\kappa|k}_i\circ(\mu\circ\operatorname{id}_T)$, où $\unsure{\kappa|k}_i$ est un
hom lacet $T\longrightarrow\pi$, et $\mu\in\mathbb N^{*}$

c) $\lambda\vert I$ n'est pas l'hom trivial

d) la section envisagée \inserted{$E_{\widetilde f}$ de $E_{U_{\widetilde K}}$ sur $E_{\widetilde K}$} ne \deleted{\unsure{se déduit|se déduise|question}}\inserted{provient} pas
d'une section de $E_{U_{\widetilde O}}$ sur $E_{\widetilde O}$

e) La section rationnelle \deleted{régulière} $\widetilde f$ de $\widehat U$ sur
$\widetilde S$ n'est pas régulière (i.e. la sect. rat.
$f$ de $U$ sur $S$ n'est pas définie en $s$)

\deleted{f)} \deleted{Si}\inserted{Désignons par} $\widetilde f$ \deleted{\unsure{désigne|désignons} \unsure{aussi}} la section de $X_{\widetilde O}$ sur $\widetilde S$
définie par $\widetilde f$ \unsure{\deleted{alors}|loc.|\deleted{les}},
\deleted{$\begin{array}[b]{@{}c@{\ }l@{}}
&\inserted{\unsure{i\in I(k)}\,\unsure{\simeq|\sim}:}\\
\exists&i\ \text{\unsure{tel que}}
\end{array}$} On a
\[
\widetilde f(s)\in I=X_s\setminus U_s.
\]
De plus, dans ce cas le $i$ dans b) n'est
autre que $\widetilde f(s)$.

Nous plaçant dans ce cas, \uline{\unsure{ns|on|nous} pouvons donc
identifier le groupe de décomposition} $N_{s'}\simeq E_{\widetilde K}$
à un sous-groupe de $\deleted{\unsure{\pi_s}}E_{X_s}$, \inserted{ext. de $E_{k(s)}$ par $\pi$,} qui est
d'indice fini $\mu$ dans \inserted{un} $N(L_i)$. Ainsi,
notre scindage \inserted{$E_{U_K}$ sur $E_K$, ou} de $E_{U_{\widetilde K}}$ \inserted{$E_{X_s}$} sur $E_{\widetilde K}$,
qui ne provient pas d'un scindage de
\margin{89}

\scanpage{91}
\mspage{45}
$\mathcal E_U$ sur $\mathcal E_S$ ni $=$ de $\mathcal E_{\widetilde U}$ sur $\mathcal E_{\widetilde{\mathcal O}}$, définit néanmoins, \deleted{\unsure{leur|le|sur}} \unsure{le|de} sous-groupe $\mathcal E_{\widetilde K}=N_{s'}$ de $\mathcal E_K$ (qui est extension d'un $\mathcal E_k$ par un facteur trivial kummérien $T\underset{\deleted{\unsure{de|d}}}{\simeq}I_{s'}$, où $k$ est un corps de degré de \unsure{tr.-scend|tr.-scendance} $\deg\operatorname{tr}K-1$, lorsque $S$ est de type fini sur un corps de \unsure{réf|$N$}, disons $\mathbf Q$) \uline{presque} un relèvement --- seul le facteur $I_{s'}\simeq T$ y obstrue encore : en tous cas $N_{s'}\simeq\mathcal E_{\widetilde K}$ est plongé dans un $N_{\mathcal E_{X_s}}(L_i)$, et toute section de cette extension de $E_{k(s)}$ par $T$ donne donc une section de $\mathcal E_{X_s}$ sur $\mathcal E_{k(s)}$ --- mais de \uline{deuxième} espèce.

Une \deleted{\unsure{f}} autre façon de dire ces choses est la suivante. Une section rationnelle $f$ de $\frac{U}{S}$ sur $S$, de domaine de définition l'ouvert $S'$ de $S$, \deleted{\unsure{peut|permet}} peut s'interpréter en termes de la famille des pts $f(s)\in U_s$, où $s$ parcourt l'ens des pts de $S$
\margin{90}

\scanpage{92}
-- ou seulement, \unsure{sans} doute, l'ens des pts fermés (si $S$ est de type fini sur un corps, tel $\mathbb Q$ disons). En termes de sections \unsure{de extension|d'extension}, on peut dire que la section $\mathcal E_{f_K}$ de $\mathcal E_{U_K}$ sur $\mathcal E_K$ donne naissance à une famille de sections des $\mathcal E_{X_s}$ sur les $\mathcal E_{k(s)}$ pour $s\in S'$. On peut \deleted{\unsure{interpréter}} \inserted{préciser} ainsi cette correspondance \unsure{entre|dentre} sections de $\mathcal E_{U_K}$ sur $\mathcal E_K$, et sections des $\mathcal E_{X_s}$ sur les $\mathcal E_{k(s)}$, en généralisant ce qu'on a vu dans le cas où $s$ est de codim $1$ : On introduit le normalisé $S'$ de $S$ dans $\overline K$, la fibre en $s$ s'identifie à l'ens des idéaux maximaux de la clôture normale de $\underline{\mathcal O}_{S,s}$ dans $\overline K$, et si $N_{s'}$ désigne le \deleted{\unsure{normalisateur}} \inserted{stabilisateur} dans $\mathcal E_{\overline K}$ d'un $s'$ de $S'_s$, on a encore une structure d'extension
\[
(29)\qquad
1\mathrel{\text{---}}I_{s'}\longrightarrow N_{s'}\longrightarrow {\deleted{\unsure{\operatorname{Gal}}}\inserted{\mathcal E}}_{k(s)}\longrightarrow 1
\]
\margin{91}

\scanpage{93}
\mspage{46}
où $I_{s'}$ \deleted{est le groupe} \inserted{$\simeq\mathcal E_{\widetilde K}$}, $\widetilde{\widetilde K}$ étant le corps
des fractions du hensélisé (\deleted{\inserted{\unsure{ou}}}\deleted{\unsure{pas|non}} strict)
de $\underline{\mathcal O}_{S,s}$. ($N_{s'}$ \deleted{est celui \unsure{du|qui} \unsure{correspond|corps de|complété}} \inserted{s'identifie de $=$ à}
$\mathcal E_{\widetilde K}$, où $\widetilde K$ correspond au hensélisé pas strict)

Ceci posé, on trouve comme précédemment,
si le composé \deleted{de} $\mathcal E_K\mathrel{\text{---}}\mathcal E_{U_K}\mathrel{\text{---}}\mathcal E_U$
s'annule sur $I_{s'}$ \inserted{(a priori, il envoie $I_{s'}$ dans $\unsure{\pi|T}$)} \deleted{\unsure{pour}} (ce qui signifie
justement \unsure{\deleted{que}|que} $s'\in S'\ldots$) que un hom
$\mathcal E_{k(s)}\longrightarrow\mathcal E_U$, ou si on préfère
$\mathcal E_{k(s)}\longrightarrow\mathcal E_{U_{\widetilde{\mathcal O}_s}}\xleftarrow{\sim}\mathcal E_{X_s}$, qui \deleted{peut} \inserted{donc} s'inter%
\unsure{prete|prète|préter} comme une section de $\mathcal E_{X_s}$ sur
$\mathcal E_{k(s)}$. Et que se passe-t-il pour
$s\notin S'$ ? \deleted{\unsure{Ainsi|Aussi}} On a vu que
$D=S\setminus S'$ est purement de codim $1$
(c'est l'image inverse du diviseur
$T$ sur $X$ par la fonction $f$ de $X$ sur $S$,
\unsure{definie|définie} si $g\geq1$ ou si $S$ de dim $1\ldots$)
en fait le $\mu=\mu_\alpha$ de tantot, relatif à
\margin{92}

\scanpage{94}
une composante $D_\alpha=\overline{\{s_\alpha\}}$ de ce diviseur, n'est autre que sa \uline{multiplicité} en $D_\alpha$.

Ceci rappelé, on se demande donc de \unsure{ce|en} qui se passe pour les $s\in\bigcup D_\alpha$, \deleted{\unsure{On|(les}} \inserted{et} \deleted{voit que} \inserted{\unsure{regarder.|regardons}} ceux qui sont $\in D_\alpha$, $\alpha$ \unsure{pris|fixés|fixé}, i.e. qui spécialisent un $s_\alpha$ fixé. On peut dire que la situation définit une extension $\mathcal E_{U_{s_\alpha}}$ de $\mathcal E_{k(s_\alpha)}$ par $\pi$ --- déduite, comme \unsure{tantôt|tantot} pour $S$ et $K$, d'une courbe \unsure{relatif|relative} $U_{D_\alpha}$ sur $D_\alpha$ --- ou $\unsure{U_{S_\alpha}|U_{s_\alpha}}$ sur un ouvert régulier convenable de $D_\alpha$, si on préfère --- en passant au pt générique. Cette fois-ci, pour la fibre générique, on n'a plus des \deleted{\unsure{sections|section}} pt. rat. sur $k(s_\alpha)$ de $\unsure{U_{s_\alpha}|U_{S_\alpha}}$ --- on n'a qu'un pt : l'infini $i_\alpha\in X_{s_\alpha}\setminus U_{s_\alpha}$, rat sur $\unsure{k_\alpha\xi^{\cdot}|K_\alpha\xi^{\cdot}}$. Mais on dispose \inserted{maintenant} d'un $L_i\subset\pi$ relatif : cet $i_\alpha$, justement, de façon ess. canonique,
\margin{93}

\scanpage{95}
\mspage{47}
qui nous incite justement : regarder la \unsure{classes|classe} de scindages, de $\mathcal E_{X_{\unsure{B.|B_0}}}$ sur $\mathcal E_{k(s)}$ relatifs à cet $L_i$ --- et notamment ceux relatifs au choix \unsure{d'un|d'une} uniformisant de $\mathcal O_{S,s_\alpha}$, i.e. d'une équation locale de $D_\alpha$ au vois. de son pt gén. $s_\alpha$, \deleted{\unsure{(plus gén.)}} \inserted{(qui \unsure{intervient|interviennent} \unsure{p\textsuperscript{r}|par} le \unsure{brin|biais|bon|bout} qu'elles définissent $\mathfrak m_{s_\alpha}/\mathfrak m_{s_\alpha}^{2}$)} Quant à ces sections-là, on voit comment elles se spécialisent, du moins c'est clair en les pts de $D_\alpha$ \unsure{sur|en} les quels la section rationnelle envisagée du faisceau conormal $\underline N_{D_\alpha,S}^{\unsure{r|+|\vee}}$ (si on suppose $D_\alpha$ régulier, \unsure{\deleted{)}|)|\&} \unsure{de|donc|du} \unsure{..|:} $\underline N_{D_\alpha/S}$ inversible) reste régulière, et inversible.

On peut se proposer d'étudier \deleted{comment} \inserted{ce qui} se passe en les autres points --- i.e. la spécialisation des sections \inserted{de deuxième type} d'extensions de groupes profinis, en faisant une étude analogue à la précédente, pour cette question, Mais cela me semble \deleted{\unsure{unique|moins}} \inserted{d'un} intérêt moins brulant --- car cela \unsure{vient|n'est}
\margin{94}

\scanpage{96}
plus une question pour, des courbes relatives
\uline{globales}, mais pour des ``éléments de
courbe'' relatifs --- savoir, pour le
complété formel d'une courbe relative
$X/B$ le long d'une section $\unsure{g_i|g}$,
avec le diviseur de cette section \unsure{distingu|disting} % word meets the right paper edge; no ending supplied
dans la structure.

Je vais plutôt revenir à la situation
où on part de une section de première
espèce d'un $\mathcal E_{U_K}$ sur $\mathcal E_K$, associée
donc à un pt \unsure{rationnel|rationnelle} $f_K$ de $U_K$ sur $K$.
On peut donc dire (indépendamment de
la donnée du modèle $S$ de $K$ sur
lequel $U_K$ \unsure{pouvait|pourrait} déjà se définir
en \unsure{tant|\deleted{tant}} que schéma relatif élémen% source line-break hyphen
taire) que la dite section de
$\mathcal E_{U_K}$ sur $\mathcal E_K$ \inserted{$(29)\quad\unsure{(f_K)_*|\mathcal E f_K|f_K}=\psi:\mathcal E_K\longrightarrow\mathcal E_{U_K}$} définit, pour \unsure{tt|tte} section
de $\mathcal E_K$ sur $\mathcal E_{\mathbf Q}=\Gamma_{\mathbf Q}$ ou sur un sous-%
groupe ouvert \inserted{$\Gamma'$} de celui-ci, une section
\margin{95}

\scanpage{97}
\mspage{48}
correspondante de $\mathcal E_{U_K}$ sur $\Gamma'$ --- d'où, en particulier, des actions de $\Gamma'$ sur $\pi$.

(\={M} chose d'ailleurs pour des sous-corps $k'$ quelconques de $K$, quand on regarde $\mathcal E_K$ comme étant sur $\mathcal E_{k'}$ : la section (29) donne une famille de sections de $\mathcal E_{U_K}$ sur des $\mathcal E_k$, \inserted{\uline{ou} des sous-groupes ouverts $\Gamma'$ de ceux-ci,} en la composant. avec des sections de $\mathcal E_K$ sur des $\mathcal E_{k'}$ --- des sous-groupes $\Gamma'\subset\mathcal E_k$. J'ai envie d'avoir une vue d'ensemble sur cette famille de sections, et de la mettre en relation avec la famille de sections des $\mathcal E_{X_s}$ sur les $\mathcal E_{k(s)}$ \inserted{définie par (29)}, définissant des actions effectives (pas seulement ext) des $\mathcal E_{k(s)}$ sur $\pi$. Pour ces derniers, les $k(s)$ n'apparaissent pas comme des sous-corps de $K$, mais plutôt comme des $\underline{\mathcal O}_s/\mathfrak m_s$ --- des corps résiduels de sous-anneaux : Cependant, si on suppose $S$ de dim \deleted{\unsure{finie}} \inserted{relative 1} sur
\margin{96}

\scanpage{98}
\margin{On aura intérêt ici de prendre pour $S$ le mod\unsure{è|e}le complet régulier can. de $K$ sur $k\ldots$}
un corps $k$, alors les $k(s)$ (pour $s\ne\eta$, pt gén.) apparaissent comme des extensions finies de ce $k$, et tout $s\in S$ définit une classe de conjugaison de sous groupes $\unsure{\Delta_{s'}|I_{s'}}^{\unsure{\deleted{\xi}|\xi}}$ dans $\unsure{\mathcal E|K}_K$ (pour sa structure : \unsure{celle|cette|la $M$} $=$ \inserted{\unsure{extension|intérieure|extérieure}} relative sur $\mathcal E_k\ldots$), \inserted{et les $N_{s'}$ qui} \deleted{\unsure{ce|qui}} sont justement les sous groupes \unsure{de|dé} \deleted{\unsure{normalisateurs}} \deleted{\unsure{inertie}} \inserted{décomposition} \unsure{de|des} tantôt, et on peut regarder les sections de l'extension \inserted{de $N_{s'}$ sur} des $\mathcal E_{k(s)}\underset{\inserted{\text{ouvert}}}{\subset}\mathcal E_K$ sur (sections donc ``de deuxième espèce'') et les composer avec la section donnée. Dans le cas où \deleted{$f$} \inserted{la section rat. $f$ de $U_K$ sur $S$} est \uline{définie} en $s$, il résulte aussitôt de ce qui précède que la section \inserted{partielle} obtenue de $\mathcal E_{U_K}$ au dessus de $\mathcal E_{k(s)}\hookrightarrow\mathcal E_K$. est aussi de première espèce, et\unsure{.} \deleted{elle \unsure{est}} (en tant qu'op. \deleted{\unsure{ext|est}}. \inserted{effective} sur $\pi$) \deleted{\unsure{est}} n'est autre que celle qui est définie par le pt $f(s)\in U_s(k(s))$. Dans
\margin{97}

\scanpage{100}
dégénérée. On peut passer \unsure{encore|meme} à
l'hensélisé (non strict) de $\underline{\mathcal O}_s$, \unsure{on|en}
trouve la situation d'un trait hensélien
$S$, à corps résiduel $k$ (ici $\mathbf Q$ --- ou une
extension finie de $\mathbf Q$), d'une
courbe $U_\eta$ de type $(g,\nu)$ (ici, $(1,1)$)
sur son pt générique $\eta$, on s'intéresse
à l'extension $\mathcal E_{U_\eta}$ de $\mathcal E_\eta$ par $\pi$,
où ici $\mathcal E_\eta$ s'inscrit encore dans
\[
(30)\qquad 1\longrightarrow T\unsure{\longrightarrow|-}\mathcal E_\eta\longrightarrow\mathcal E_k\unsure{\longrightarrow|-}1
\]
\deleted{On est ici \unsure{en car. $0$|encore}}
\inserted{en plus de}
\[
(31)\qquad \unsure{..|:}\,1-\pi-\mathcal E_{U_\eta}\longrightarrow\mathcal E_\eta\longrightarrow1
\]
Ici l'action extérieure de $T$ sur
$\pi$ n'est \inserted{\unsure{?}} sûrement pas triviale (je
conjecture qu'elle l'est ssi il y
a ``bonne réduction'') --- donc le groupe
$\mathcal E_k$ n'opère pas lui $=$ extérieurement
sur $\pi$. Mais tout scindage de (30)
\margin{mais je l'ai
démontré, dans le
contexte des VA,
il y a belle
lurette ! Ça doit bien
être connu sur
$H_1$ sûrement \unsure{!!|!}}
\margin{99}

\scanpage{101}
\mspage{50}
définit une extension de $\mathcal E_k$ par $\pi$, donc une action extérieure --- est-elle ``admissible'' \unsure{?|\deleted{?}}, i e. définie par une courbe algébrique ? Sans doute pas --- \uline{quelle} courbe serait-ce, si on vise la courbe ``réduite'' de type $(g,\nu)$ (inexistant !) ?
\margin{mais il y a tjrs une arithmetisation'' de $\pi$ innocente !}

Dès lors, si on a une section de (30) qu'on compose avec une de (31) --- donc une action \inserted{effective} de $\mathcal E_k$ sur $\pi$ --- qu'attendre d'une telle action ?

\unsure{Sans|\deleted{Sans}} m'y attendre, me voici à tâter un peu la structure : l'\unsure{$\infty$|os} des variétés modulaires de Teichm\unsure{ù|ü}ller !

S'il faut hasarder des conjectures dans cette situation scabreuse, ce pourrait-être celle ci :

\uline{Conjecture-à-tout-hasard} : Les conditions suivantes sont équivalentes

\noindent a) $U_\eta$ a bonne réduction sur $S$

\noindent b) L'action extérieure de $T$ sur $\pi$ est triviale
\margin{100}

\scanpage{102}
Ce qui signifie aussi que tt \inserted{(ou un)} scindage
de 31 --- p.ex. défini par un pt de $U_\eta$
\unsure{rat|rt}/$K$ --- induit un hom $T\longrightarrow\pi$)

\noindent c) L'action de $T$ sur $\pi_{\mathrm{ab}}=H_1(U_{\bar\eta})$ est
triviale.

\noindent d) Item pour $H_{\unsure{1|n}}(U_{\bar\eta},T_\ell)$, si $\ell$ est premier
à la car. résiduelle. (\uline{NB} \deleted{\unsure{Pour car. $p>0$|Pour $\ell\ne p$|Pour $\ell=p$}} \inserted{$\ell=p>0$}\unsure{,|)}
on doit pouvoir l'exprimer par un critère
sur le groupe $p$-divisible\ldots)

\noindent e) (En termes d'une section de (30),
définissant une extension de $\mathcal E_k$ par $\pi$,
et supposant maintenant $k$ de type
fini sur $\mathbf Q$) La dite extension est
admissible (quitte à passer à un
sous groupe ouvert \inserted{de $\mathcal E_k$,} au besoin), ie.
peut être défini par une courbe alg.
sur \inserted{$k$, ou sur} une extension finie de $k$.

\noindent f) (En termes d'une section de (30)
et d'une de (31)., permettant de
faire opérer $\mathcal E_k$ sur $\pi$) : Cette
action \inserted{n'}est \inserted{pas trop} loufoque (à définir convenablement !)% The end of convenablement continues on a short lower baseline; no hyphen is written.
\margin{101}
